% Part: proof-theory
% Chapter: cut-elimination
% Section: introduction

\documentclass[../../../include/open-logic-section]{subfiles}

\begin{document}

\olfileid{pt}{cut}{int}

\olsection{جان پچھان}

\Cut{} دا قاعدہ ایہ اے:
\[
\Axiom$\Gamma \fCenter \Delta, !A$
\Axiom$!A, \Pi \fCenter \Lambda$
\RightLabel{\Cut}
\BinaryInf$\Gamma, \Pi \fCenter \Delta, \Lambda$
\DisplayProof
\]
(!!{formula}~$!A$ نوں \emph{قطع فارمولا} آکھدے نیں۔)

\Cut{} گینٹسن دے اصلی سیکوئنٹ حساب~\Log{LK} وچ شامل سی۔ گینٹسن دا \emph{Hauptsatz} («مُکھ قضیہ») کہندا اے کہ ہر سیکوئنٹ جیہڑا~\Log{LK} وچ !!{provable} اے، \Cut{} دا قاعدہ ورتے بغیر وی !!{provable} اے۔ یعنی \Cut{} نوں \Log{LK} دے ثبوتاں وچوں «ہٹایا» جا سکدا اے۔ ساڈے نظام \Log{G1c} تے \Log{G3c} وچ \Cut نہیں، پر اوہناں لئی وی ایہ سوال پچھیا جا سکدا اے۔ کی \Cut{} اوہناں !!{proof}s وچوں ہٹایا جا سکدا اے جیہڑے \Cut{} نال \Log{G1c} دے باقی قاعدے (یا \Log{G3c} دے قاعدے) ورتدے نیں؟ ساڈیاں قائم اصطلاحاں وچ ایہو سوال ایویں اے۔ کی \Cut{}، \Log{G1c} (یا \Log{G3c}) دا قبول جوگ قاعدہ اے؟

قطع ہٹاؤن دا قضیہ شاید ثبوتی نظریے دا سبھ توں اہم تے مشہور نتیجہ اے۔ پہلاں ایہ قائم کردا اے کہ \Log{G1c} تے \Log{G3c} نوں اک ایہو جیہے قاعدے دی لوڑ نہیں جیہڑا پہلی نظرے صاف ضروری لگدا اے۔ جدوں تسی اصل ثبوتاں نوں رسمی شکل دین دی کوشش کرو، تاں دِسدا اے کہ معاون قضیے ورتن دا ڈھنگ چاہیدا اے۔ ریاضی دان نتیجے ثابت کردے نیں تے فیر ہور نتیجیاں دے ثبوتاں وچ اوہناں نوں ورتدے نیں۔ سو پہلاں نتیجے $L$ (معاون قضیے) دا ثبوت، سیکوئنٹ $\Sequent L$ دے !!{proof} وجوں، رسمی شکل وچ لکھ سکدے او۔ فیر دوجے نتیجے~$R$ دا ثبوت، جیہڑا~$L$ ورتدا اے، سیکوئنٹ $L \Sequent
R$ دے !!{proof} وجوں لکھدے او۔ تہاڈے کول~$R$ دا ثبوت اے، یعنی بس !!a{proof} اے $\Sequent R$ دا، ایہ کِویں پتا لگے؟ دوہاں !!{proof}s نوں جوڑن دا ڈھنگ چاہیدا اے، تے \Cut{} دا قاعدہ ایہ کرن دی کھُل دیندا اے۔

اسیں دسیا سی کہ \Log{G1c} تے \Log{G3c} پورے نیں، سو اوہ ہر اوہ سیکوئنٹ ثابت کردے نیں جیہڑا ثابت ہو سکنا چاہیدا اے، یعنی سارے معتبر سیکوئنٹ۔ ایہ سدھا ثابت کیتا جا سکدا اے۔ پر جدوں گینٹسن لکھدا سی، اوہدوں بس مسلّماتی اخذ دے نظاماں بارے پتا سی کہ اوہ پورے نیں۔ \Log{LK} دی تمامیت وکھاؤن لئی گینٹسن نے مسلّماتی نظام دے !!{derivation}s نوں \Log{LK} دے !!{proof}s وچ بدلّن دا طریقہ دِتّا۔ ایس تبدیلی وچ \Cut{} ضروری طور اُتے ورتیا گیا سی۔ قطع ہٹاؤنا کلاسیکی منطق لئی ای اہم نتیجہ نہیں، کئی ہور منطقاں دے وی سیکوئنٹ حساب نیں۔ کجھ منطقاں وچ اوہناں دے سیکوئنٹ حساباں دی تمامیت دا سدھا ثبوت اوکھا یا ناممکن اے۔ سو دوجے نظاماں توں تبدیلی (جیہڑی \Cut ورتدی اے) ای تمامیت وکھاؤن دا اکّو ڈھنگ اے۔ فیر ثبوتی نظریے نال قطع ہٹاؤنا ثابت کرنا ضروری اے تاں جو قائم ہووے کہ \Cut{} وادھو اے تے \Cut{} توں بناں نظام پورا اے۔

قطع ہٹاؤنا ایس لئی وی اہم اے کہ ایہنوں ہور، اپنے آپ وچ دلچسپ، نتیجے کڈھن لئی ورتیا جا سکدا اے۔ اسیں دو صورتاں ویکھاں گے، کریگ دا درمیانی فارمولا والا معاون قضیہ تے ہیربرانڈ دا قضیہ۔

قطع ہٹاؤنا ثابت کرن دے خیالاں وچ عام طور اُتے ایہ وکھایا جاندا اے کہ \Cut{} ورتن والا !!a{proof} اوس توں بناں ثبوت وچ کِویں بدلیا جا سکدا اے۔ ایہ بدلاواں عموماً «مقامی» ہوندے نیں۔ یعنی اسیں !!
a{proof} دا اک حصہ لیندے آں (عام طور اُتے \Cut{} استنباط اُتے مُکدا ذیلی !!{proof}) تے اوہدی تھاں اوسے سیکوئنٹ دا ہور ذیلی !!{proof} رکھدے آں، جیہدے وچ \Cut{} استنباط ہٹا دِتّا گیا اے یا اوہدی تھاں ہور استنباط رکھے گئے نیں۔ ایہو جیہیاں دلیلاں قدم وار الگورتھم دیندیاں نیں جیہڑے \Cut{} والے !!{proof}s نوں اوس توں بناں ثبوتاں وچ بدل دیندے نیں۔ قطع ہٹاؤن دے طریقے ورتن ویلے ایس توں ہور جانکاری ملدی اے۔ مثلاً درمیانی فارمولے دے معاون قضیے دے نمونیاں دے نظریے والے ایہو جیہے ثبوت نیں کہ جے «$!A \lif !B$» معتبر اے، تاں اک درمیانی فارمولا موجود اے (ہُن لئی ایہ چھڈو کہ درمیانی فارمولا کیہ اے)۔ قطع ہٹاؤن والی ثبوتی دلیل اک الگورتھم دیندی اے جیہڑا !!a{proof} لے کے $!A
\lif !B$ دا درمیانی فارمولا موڑدا اے۔ نمونیاں دے نظریے والی دلیل دے مقابلے وچ، ثبوتی دلیل \emph{تعمیری} اے۔

قطع ہٹاؤن دا قضیہ ثابت کرن دے دو مشہور ڈھنگ نیں۔ اک، جیہڑا گینٹسن توں آؤندا اے، وکھاؤندا اے کہ ثبوت وچوں سبھ توں اُتلے قطع ہٹا سکدے آں، تے فیر اوہ اُتلے قطع ہٹائی جاندا اے جدوں تک کوئی نہیں رہندا۔ دوجا، جیہڑا مُڈھوں شوٹے تے ٹیٹ دا اے، !!a{proof} دے سبھ توں گنجل دار قطعاں اُتے دھیان دیندا اے تے وار واری اوہ قطع ہٹاؤندا اے جیہڑے سبھ توں اُچّے درجے دے نیں، ایس معنی وچ کہ ہور کسے قطع دا قطع !!{formula} اوہناں توں ڈونگھا نہیں۔ دوہاں دے اپنے فائدے تے اوکڑاں نیں۔ کجھ نظاماں وچ اک ڈھنگ کم کردا اے تے دوجا لاؤنا اوکھا یا ناممکن اے۔ سو اسیں دوویں ڈھنگ دساں گے۔ اسیں~\Log{G3c} لئی قطع ہٹاؤنا ثابت کراں گے۔ اصل وچ اسیں \Cut{} ہٹاؤن دی تھاں \emph{سانجھے سیاق والا قطع}~\CutCS ہٹاؤنا وکھاواں گے:
\[
\Axiom$\Gamma \fCenter \Delta, !A$
\Axiom$!A, \Gamma \fCenter \Delta$
\RightLabel{\Cut}
\BinaryInf$\Gamma \fCenter \Delta$
\DisplayProof
\]
جے سیکوئنٹ حساب کمزور کرن تے سکیڑن دی کھُل دیندا اے (بھاویں \Log{G1c} وانگ اصلی قاعدیاں وجوں یا \Log{G3c} وانگ قبول جوگ قاعدیاں وجوں)، تاں \Cut{}، \CutCS دا کم کر سکدا اے، تے اُلٹا وی۔ اسیں \CutCS{} چُندے آں \Cut دی تھاں، کیوں جے پہلا طریقہ~\CutCS لئی سمجھاؤنا سوکھا اے تے دوجا طریقہ بس~\CutCS لئی کم کردا اے۔

\end{document}
